% Generated by code/fill_numbers.py from chuong3_credit_b.tpl -- edit the template, not this file.
\subsection{Kết quả}\label{sec:credit-results}

\begin{table}[t]
\centering
\caption[Kết quả trên dữ liệu vỡ nợ thẻ tín dụng]{Kết quả trên dữ liệu vỡ nợ thẻ tín dụng (trung bình $\pm$ độ lệch chuẩn qua $5$ lần chia $70/30$). AUC: diện tích dưới đường ROC; ĐCX-CB: độ chính xác cân bằng; F1 tính cho lớp vỡ nợ; ngưỡng được chọn trên tập huấn luyện sao cho tỉ lệ dự báo vỡ nợ bằng tỉ lệ vỡ nợ quan sát. Cột cuối: số thực cần lưu (trung vị); thời gian dự đoán được đo riêng ở \cref{sec:cost}.}
\label{tab:credit}
\footnotesize
\setlength{\tabcolsep}{3pt}
\renewcommand{\arraystretch}{1.18}
\begin{tabular}{@{}l c c c r@{}}
\toprule
\textbf{Mô hình} & \textbf{AUC} & \textbf{ĐCX-CB (\%)} & \textbf{F1 (\%)} & \textbf{Số thực}\\
\midrule
Hồi quy logistic & $0{,}724\pm0{,}009$ & $68{,}7\pm0{,}6$ & $51{,}2\pm0{,}9$ & 79\\
Thẻ điểm (logistic, biến rời rạc hóa) & $0{,}773\pm0{,}006$ & $70{,}3\pm0{,}6$ & $53{,}7\pm1{,}0$ & 381\\
SVM tuyến tính & $0{,}706\pm0{,}008$ & $68{,}5\pm0{,}8$ & $50{,}9\pm1{,}4$ & 79\\
\midrule
PWL-LS-SVM, cộng tính, lưới đều & $0{,}772\pm0{,}006$ & $70{,}3\pm0{,}5$ & $53{,}6\pm0{,}9$ & 547\\
PWL-LS-SVM, cộng tính, phân vị & $0{,}776\pm0{,}006$ & $70{,}4\pm0{,}7$ & $53{,}8\pm1{,}1$ & 349\\
\quad và ràng buộc đơn điệu (Mục~\ref{sec:monotone}) & $0{,}775\pm0{,}006$ & $70{,}3\pm0{,}8$ & $53{,}7\pm1{,}2$ & 349\\
PWL-C-SVM, cộng tính, phân vị & $0{,}744\pm0{,}007$ & $67{,}8\pm0{,}8$ & $49{,}9\pm1{,}4$ & 349\\
PWL-LS-SVM, HH & $0{,}743\pm0{,}007$ & $69{,}5\pm0{,}6$ & $52{,}4\pm1{,}0$ & 6.631\\
\midrule
SVM-RBF & $0{,}716\pm0{,}010$ & $69{,}2\pm0{,}6$ & $52{,}0\pm1{,}0$ & 243.134\\
MLP-ReLU & $0{,}768\pm0{,}005$ & $70{,}0\pm0{,}4$ & $53{,}3\pm0{,}7$ & 6.605\\
Tổ hợp cây tăng cường & $0{,}779\pm0{,}006$ & $70{,}3\pm0{,}7$ & $53{,}7\pm1{,}1$ & 4.636\\
\bottomrule
\end{tabular}
\end{table}


\begin{figure}[t]
\centering
\includegraphics[width=0.62\linewidth]{fig_credit_roc.pdf}
\caption[Đường cong ROC trên dữ liệu tín dụng]{Đường cong ROC trên tập kiểm tra của lần chia đầu tiên cho bốn mô hình; con số trong ngoặc là AUC trung bình qua năm lần chia.}
\label{fig:credit-roc}
\end{figure}

\Cref{tab:credit} và \cref{fig:credit-roc} cho thấy các mô hình chia thành hai nhóm rõ rệt. Nhóm dẫn đầu gồm tổ hợp cây tăng cường (AUC trung bình 0{,}779), PWL-LS-SVM cộng tính với nút phân vị (0{,}776), thẻ điểm logistic trên các biến đã rời rạc hóa (0{,}773) và PWL-LS-SVM cộng tính với lưới đều (0{,}772). Nhóm còn lại gồm hồi quy logistic trên các biến gốc (0{,}724), SVM tuyến tính (0{,}706) và SVM hạt nhân Gauss (0{,}716); mạng nơ-ron (0{,}768) đứng giữa hai nhóm. Như vậy, AUC của PWL-LS-SVM cộng tính với nút phân vị chỉ kém tổ hợp cây tăng cường một chút. Độ chính xác cân bằng và độ đo F1 ở ngưỡng đã chọn dẫn tới cùng kết luận: bốn mô hình của nhóm dẫn đầu gần như ngang nhau và cao hơn rõ rệt hồi quy logistic, SVM tuyến tính và SVM hạt nhân Gauss. Kết quả rất ổn định giữa các lần chia: độ lệch chuẩn của AUC chỉ khoảng 0{,}006.

Khoảng cách 0{,}052 đơn vị AUC giữa PWL-LS-SVM và hồi quy logistic trên các biến gốc cần được hiểu đúng. Phần lớn khoảng cách đó đến từ quan hệ phi tuyến giữa các trạng thái trả nợ và rủi ro (\cref{fig:credit-overview}a), mà một mô hình tuyến tính theo PAY\_0 không mô tả được. Trong thực tế, các thẻ điểm không dùng trực tiếp biến gốc mà rời rạc hóa từng biến thành các khoảng rồi mới hồi quy logistic~\cite{thomas2002}, và thẻ điểm xây dựng theo cách đó ở \cref{tab:credit} đã thu hẹp phần lớn khoảng cách. Thẻ điểm cũng là đối chứng công bằng nhất cho PWL-SVM cộng tính: cả hai đều là mô hình cộng tính, đều dựa trên các khoảng của từng biến, chỉ khác ở chỗ đóng góp của mỗi biến là hàm bậc thang trong thẻ điểm và là đường gấp khúc liên tục trong PWL-SVM. So với thẻ điểm, PWL-LS-SVM cao hơn ở bốn trong năm lần chia, với chênh lệch trung bình 0{,}002 đơn vị AUC, trong khi hai mô hình có kích thước tương đương (349 và 381 số thực). Chênh lệch này nhỏ, nhưng trong chấm điểm tín dụng, ngay cả những cải thiện nhỏ về khả năng phân biệt cũng có thể mang lại lợi ích kinh tế đáng kể khi áp dụng cho một danh mục lớn~\cite{lessmann2015}.

Ba quan sát khác đáng chú ý. Thứ nhất, lưới phân vị tốt hơn lưới đều (0{,}776 so với 0{,}772) với mô hình nhỏ hơn (349 so với 547 số thực), đúng như phân tích ở \cref{fig:credit-overview}b. Thứ hai, với cùng ánh xạ, PWL-C-SVM cho AUC thấp hơn hẳn PWL-LS-SVM (0{,}744). Điều này có thể giải thích bằng hàm mất mát: mất mát bản lề chỉ quan tâm đến các điểm nằm gần hoặc vi phạm lề, nên giá trị hàm quyết định ở xa biên không mang nhiều thông tin về mức độ rủi ro; còn mất mát bình phương của LS-SVM, tương đương hồi quy trên nhãn $\pm1$ (\cref{prop:ridge}), buộc điểm số của mọi khách hàng phải phản ánh nhãn của họ, và vì thế cho một thứ hạng tốt hơn. SVM hạt nhân Gauss, cũng dùng mất mát bản lề, gặp cùng hạn chế. Thứ ba, ánh xạ HH (0{,}743) kém ánh xạ cộng tính: trên dữ liệu này, ảnh hưởng của từng biến quan trọng hơn tương tác giữa các biến, và các siêu phẳng ngẫu nhiên lãng phí khả năng biểu diễn, phù hợp với kết luận của Mục~\ref{sec:sensitivity}.

\subsection{Diễn giải mô hình}\label{sec:credit-interp}

Vì ánh xạ là cộng tính, hàm quyết định của PWL-LS-SVM tách thành tổng các đóng góp của từng biến,
\begin{equation}\label{eq:additive-decomp}
f(x)=w_0+\sum_{i=1}^{n}h_i\bigl(x(i)\bigr),
\end{equation}
trong đó mỗi $h_i$ là một hàm tuyến tính từng phần một biến, có không quá chín điểm gãy đặt tại các phân vị của biến đó (\cref{prop:1d}). Mỗi $h_i$ có thể được vẽ ra và đọc trực tiếp, giống như một thẻ điểm trong đó số điểm cộng thêm cho mỗi khách hàng là một đường gấp khúc theo giá trị của biến. \Cref{fig:credit-shapes} vẽ sáu hàm $h_i$ có ảnh hưởng lớn nhất của mô hình huấn luyện trên lần chia đầu tiên, sau khi trừ đi giá trị trung bình trên tập huấn luyện; giá trị dương nghĩa là rủi ro cao hơn một khách hàng trung bình.

\begin{figure}[t]
\centering
\includegraphics[width=\linewidth]{fig_credit_shapes.pdf}
\caption[Các hàm đóng góp của mô hình PWL-SVM cộng tính]{Sáu hàm đóng góp $h_i$ có ảnh hưởng lớn nhất của PWL-LS-SVM cộng tính (nút phân vị) trên dữ liệu tín dụng, đã trừ giá trị trung bình. Nét đậm: mô hình của lần chia đầu tiên, nét chấm dọc là vị trí các nút của nó; nét mảnh: mô hình của bốn lần chia còn lại. Giá trị dương ứng với rủi ro vỡ nợ cao hơn.}
\label{fig:credit-shapes}
\end{figure}

Hàm đóng góp của PAY\_0 thể hiện rõ cấu trúc mà hồi quy logistic trên biến gốc bỏ lỡ (so sánh với \cref{fig:credit-overview}a). Ba mã không chậm trả ($-2$, $-1$ và $0$) có đóng góp thấp và gần nhau; mã $-1$ cao hơn hai mã còn lại một chút, đúng như tỉ lệ vỡ nợ quan sát được (16{,}8\% so với 13{,}2\% và 12{,}8\%). Đóng góp tăng mạnh khi khách hàng chậm trả một tháng (tỉ lệ vỡ nợ 33{,}9\%) và đạt cực đại khi chậm hai tháng (69{,}1\%); chênh lệch giữa mã $0$ và mã $2$ là $0{,}86$ đơn vị điểm, lớn hơn toàn bộ biên độ của bất kỳ biến nào khác. Đóng góp của hạn mức tín dụng giảm dần theo hạn mức --- điều hợp lý, vì hạn mức do chính ngân hàng xác định dựa trên đánh giá tín nhiệm --- và khoảng 60\% mức giảm diễn ra ở vùng hạn mức dưới $100$ nghìn Đài tệ. Trạng thái trả nợ của các tháng trước (PAY\_2, PAY\_3, PAY\_4) có hình dạng tương tự PAY\_0 nhưng biên độ nhỏ hơn nhiều: thông tin gần nhất là quan trọng nhất.

Hai đặc điểm khác của các hàm đóng góp cần được đọc thận trọng. Thứ nhất, sau mức chậm hai tháng, đóng góp của PAY\_0 và nhất là của PAY\_4 lại giảm, trái với trực giác. Đây không phải là quy luật của dữ liệu mà là hệ quả của việc có quá ít dữ liệu: chỉ 1{,}5\% khách hàng có PAY\_0 từ $3$ trở lên, tỉ lệ vỡ nợ của nhóm này (71{,}9\%) không thấp hơn của nhóm chậm hai tháng, và những khách hàng này thường chậm trả nhiều tháng liền, nên mô hình có thể phân bổ rủi ro giữa các biến PAY theo nhiều cách gần tương đương. Hệ số góc của đoạn cuối, được ước lượng từ rất ít điểm, vì vậy không đáng tin cậy; Mục~\ref{sec:monotone} sẽ khắc phục hiện tượng này. Thứ hai, các biến dư nợ sao kê BILL\_AMT1--6 tương quan rất mạnh với nhau (hệ số tương quan giữa hai tháng liền kề từ $0{,}92$ đến $0{,}95$), nên mô hình có thể chuyển ảnh hưởng qua lại giữa chúng: chẳng hạn, đóng góp của BILL\_AMT3 tăng theo dư nợ, còn đóng góp của BILL\_AMT4 lại giảm. Với những nhóm biến như vậy, chỉ tổng đóng góp của cả nhóm mới nên được diễn giải. Đây là hạn chế chung của mọi mô hình cộng tính khi các biến tương quan mạnh, không riêng gì PWL-SVM.

Các nét mảnh ở \cref{fig:credit-shapes} cho thấy các hàm đóng góp thay đổi thế nào khi dữ liệu huấn luyện thay đổi. Ở vùng có nhiều dữ liệu, hình dạng gần như không đổi giữa các lần chia: giữa phân vị $5\%$ và $95\%$ của hạn mức tín dụng, đóng góp của LIMIT\_BAL ở bốn lần chia còn lại lệch khỏi lần chia đầu tiên không quá 15\% biên độ của nó, và ở các trạng thái từ $-2$ đến $2$, đóng góp của PAY\_0 lệch không quá 10\% biên độ. Thứ hạng của các biến cũng ổn định: ở cả năm lần chia, PAY\_0 là biến có ảnh hưởng lớn nhất và LIMIT\_BAL đứng thứ hai; thứ tự của các biến tiếp theo, vốn có mức ảnh hưởng gần nhau, thì thay đổi giữa các lần chia. Ngược lại, ở các trạng thái chậm trả từ ba tháng trở lên, các lần chia cho những đường rất khác nhau: độ lệch chuẩn giữa các lần chia của đóng góp, tính trung bình trên sáu biến PAY, ở những trạng thái này lớn gấp khoảng 8 lần so với ở các trạng thái từ $-2$ đến $2$ --- thêm một bằng chứng rằng các đoạn đó không đáng tin cậy.

\begin{figure}[t]
\centering
\includegraphics[width=\linewidth]{fig_credit_explain.pdf}
\caption[Mức độ ảnh hưởng của các biến và giải thích một hồ sơ]{(a) Mức độ ảnh hưởng của mười biến quan trọng nhất, đo bằng độ lệch chuẩn của đóng góp $h_i$ trên tập huấn luyện; với các biến định danh, đó là độ lệch chuẩn của tổng đóng góp của các biến chỉ báo tương ứng. (b) Tám đóng góp lớn nhất, so với khách hàng trung bình, của hồ sơ có điểm rủi ro cao nhất trong số các khách hàng vỡ nợ ở tập kiểm tra.}
\label{fig:credit-explain}
\end{figure}

\Cref{fig:credit-explain}a xếp hạng các biến theo mức độ ảnh hưởng. PAY\_0 vượt xa mọi biến khác; tiếp theo, với mức ảnh hưởng nhỏ hơn nhiều, là hạn mức tín dụng, các trạng thái trả nợ của những tháng trước và các số tiền sao kê; các biến nhân khẩu học (học vấn, hôn nhân, giới tính, tuổi) có ảnh hưởng không đáng kể. Vì hàm quyết định~\eqref{eq:additive-decomp} là một tổng, mỗi quyết định của mô hình được giải thích một cách chính xác chứ không phải xấp xỉ: điểm của một khách hàng bằng điểm trung bình trên tập huấn luyện cộng với tổng các đóng góp của từng biến so với khách hàng trung bình. \Cref{fig:credit-explain}b minh họa điều đó với khách hàng có điểm rủi ro cao nhất trong số những người vỡ nợ ở tập kiểm tra. Điểm của khách hàng này là $0{,}96$, trong khi điểm trung bình là $-0{,}56$ và ngưỡng phân loại là $-0{,}44$. Trong phần chênh lệch $1{,}51$ so với điểm trung bình, riêng việc chậm trả hai tháng ở tháng 9 đóng góp $0{,}71$, lịch sử chậm trả ở năm tháng trước đó đóng góp thêm $0{,}63$, và hạn mức thấp ($20.000$ Đài tệ) đóng góp $0{,}13$. Một lời giải thích như vậy có thể được trình bày trực tiếp cho cán bộ tín dụng hay cho khách hàng. Với SVM hạt nhân Gauss hay tổ hợp cây, việc giải thích một quyết định đòi hỏi những kỹ thuật bổ sung, chẳng hạn giá trị SHAP~\cite{lundberg2017}, được tính sau khi huấn luyện và phức tạp hơn nhiều so với việc đọc trực tiếp các hàm $h_i$.

\subsection{Ràng buộc đơn điệu}\label{sec:monotone}

Người làm tín dụng kỳ vọng rằng rủi ro không giảm khi số tháng chậm trả tăng, và không tăng khi hạn mức tín dụng tăng. Mô hình không ràng buộc chỉ thỏa mãn một phần kỳ vọng này: ở cả năm lần chia, ít nhất một trong sáu hàm $h_i$ của các biến PAY có một đoạn với hệ số góc âm ở vùng từ trạng thái $0$ trở đi, và hàm của LIMIT\_BAL có đoạn đi lên ở cả năm lần chia. Như đã phân tích ở trên, đó là những đoạn được ước lượng từ rất ít dữ liệu, không phải là quy luật thật. Một mô hình trái với tri thức nghiệp vụ khó được chấp nhận khi thẩm định, dù độ chính xác của nó cao; vì vậy ta đưa kỳ vọng đơn điệu vào chính bài toán huấn luyện.

Với PWL-SVM cộng tính, điều này có thể làm mà vẫn giữ tính lồi. Gọi $t_{i1}<\dots<t_{iK}$ là các nút của biến thứ $i$ và $s_{i0},\dots,s_{iK}$ là hệ số góc của $h_i$ trên $K+1$ đoạn mà các nút chia ra. Khi đó
\begin{equation}\label{eq:slope-param}
h_i(t)=\sum_{j=0}^{K}s_{ij}\,G_{ij}(t),\qquad
G_{ij}(t)=\min\bigl\{(t-t_{ij})_+,\ t_{i,j+1}-t_{ij}\bigr\},
\end{equation}
với quy ước $t_{i0}=-\infty$, $t_{i,K+1}=+\infty$ (tức $G_{i0}(t)=\min\{t,t_{i1}\}$ và $G_{iK}(t)=(t-t_{iK})_+$): $G_{ij}(t)$ là độ dài phần đoạn thứ $j$ nằm bên trái $t$. Theo \cref{prop:1d}, hệ số của $x(i)$ trong ánh xạ~\eqref{eq:feature-map} bằng $s_{i0}$ và hệ số của $\bigl(x(i)-t_{ij}\bigr)_+$ bằng $s_{ij}-s_{i,j-1}$. Vì vậy~\eqref{eq:slope-param} chỉ là một cách tham số hóa khác của cùng một mô hình, và bài toán~\eqref{eq:ridge} trở thành
\[
\min_{s,\,w_0}\ \frac12\sum_{i=1}^{n}\Bigl[s_{i0}^2+\sum_{j=1}^{K}\bigl(s_{ij}-s_{i,j-1}\bigr)^2\Bigr]+\frac\gamma2\sum_{k=1}^{N}\Bigl(y_k-w_0-\sum_{i=1}^{n}h_i\bigl(x_k(i)\bigr)\Bigr)^2 .
\]
Hàm $h_i$ không giảm khi và chỉ khi $s_{ij}\ge0$ với mọi $j$, và không tăng khi và chỉ khi $s_{ij}\le0$ với mọi $j$; muốn $h_i$ đơn điệu chỉ trên một phần của trục, chẳng hạn trên nửa đường thẳng $t\ge t_{ij_0}$, ta chỉ cần áp các bất đẳng thức đó cho những đoạn nằm trong phần ấy. Thêm các ràng buộc này, ta được một bài toán bình phương tối thiểu với ràng buộc cận trên các biến --- một bài toán lồi, giải được chính xác bằng thuật toán BVLS~\cite{stark1995} có sẵn trong SciPy~\cite{virtanen2020}. Cũng như ở \cref{rem:convex}, tri thức tiên nghiệm được đưa vào mô hình mà không phải từ bỏ tính lồi. Khi không có ràng buộc nào, lời giải trùng với lời giải của \cref{prop:ridge} (sai khác dưới $10^{-9}$), điều được dùng để kiểm tra cài đặt.

\begin{figure}[t]
\centering
\includegraphics[width=\linewidth]{fig_credit_monotone.pdf}
\caption[Hàm đóng góp khi có và không có ràng buộc đơn điệu]{Hàm đóng góp của PAY\_0, PAY\_4 và LIMIT\_BAL trong mô hình không ràng buộc (nét liền) và mô hình có ràng buộc đơn điệu (nét đứt), lần chia đầu tiên. Vùng tô xám: các trạng thái chậm trả từ ba tháng trở lên, chỉ chiếm chưa tới 2\% số khách hàng.}
\label{fig:credit-monotone}
\end{figure}

Ta ràng buộc sáu hàm của các biến PAY không giảm từ trạng thái $0$ trở đi --- ba mã $-2$, $-1$ và $0$ chỉ những tình trạng không chậm trả khác nhau, không có thứ tự tự nhiên, nên được để tự do --- và hàm của LIMIT\_BAL không tăng trên toàn miền; các biến khác, các nút và hằng số $\gamma$ đã chọn cho mô hình không ràng buộc được giữ nguyên. \Cref{fig:credit-monotone} cho thấy tác động của ràng buộc: sau mức chậm hai tháng, đóng góp của PAY\_0 và PAY\_4 giữ nguyên thay vì giảm, và đường của LIMIT\_BAL trở thành đơn điệu giảm, trong khi hình dạng ở vùng có nhiều dữ liệu hầu như không đổi. Độ chính xác hầu như không thay đổi: AUC trung bình là 0{,}775, so với 0{,}776 của mô hình không ràng buộc, và mô hình có ràng buộc vẫn cao hơn thẻ điểm (0{,}773) (\cref{tab:credit}). Như vậy, với ràng buộc đơn điệu, PWL-LS-SVM cộng tính cho một mô hình vừa chính xác vừa nhất quán với tri thức nghiệp vụ.

\subsection{Khả năng triển khai}

Mô hình PWL-LS-SVM cộng tính trên dữ liệu tín dụng gồm 349 số thực: các nút và hệ số của $20$ đường gấp khúc cho các biến số, hệ số của $6$ biến chỉ báo, hệ số chặn và hai vectơ chuẩn hóa. Để chấm điểm một hồ sơ, với mỗi biến ta chỉ cần xác định giá trị rơi vào đoạn nào giữa các nút rồi tính một hàm bậc nhất, tổng cộng vài trăm phép so sánh, nhân và cộng. Một mô hình như vậy có thể được cài đặt bằng vài dòng lệnh SQL trong kho dữ liệu của ngân hàng, bằng một bảng tính, hoặc trên thiết bị có tài nguyên rất hạn chế, mà không cần thư viện học máy nào; ràng buộc đơn điệu không làm thay đổi cấu trúc này. Việc huấn luyện lại khi có dữ liệu mới chỉ mất chưa tới một giây (\cref{tab:cost-measured}). Để so sánh, SVM hạt nhân Gauss cần lưu khoảng 243.134 số thực, và tổ hợp cây tăng cường gồm từ 34 đến 47 cây quyết định với khoảng 4.636 tham số; cả hai đều khó kiểm tra và giải thích hơn nhiều.
